\standardid{1.OA.1}
\document{Math}
\sectionid{Math K-8}
\anchor{0}
\theme{}
\domain{Operations and Algebraic Thinking}
\domainid{OA}
\grade{1}
\cluster{Represent and solve problems involving addition and
  subtraction.}
\standardnumber{1}
\letter{}
\adv{}
\mod{}
\text{Use addition and subtraction within 20 to solve word problems
  involving situations of adding to, taking from, putting together,
  taking apart, and comparing, with unknowns in all positions, e.g.,
  by using objects, drawings, and equations with a symbol for the
  unknown number to represent the problem.\footnote{See Glossary,
    Table 1.}}

%%% Local Variables: 
%%% mode: latex
%%% TeX-master: "allstandards"
%%% End: 
\standardid{1.OA.2}
\document{Math}
\sectionid{Math K-8}
\anchor{0}
\theme{}
\domain{Operations and Algebraic Thinking}
\domainid{OA}
\grade{1}
\cluster{Represent and solve problems involving addition and
  subtraction.}
\standardnumber{2}
\letter{}
\adv{}
\mod{}
\text{Solve word problems that call for addition of three whole
  numbers whose sum is less than or equal to 20, e.g., by using
  objects, drawings, and equations with a symbol for the unknown
  number to represent the problem.}

%%% Local Variables: 
%%% mode: latex
%%% TeX-master: "allstandards"
%%% End: 
\standardid{1.OA.3}
\document{Math}
\sectionid{Math K-8}
\anchor{0}
\theme{}
\domain{Operations and Algebraic Thinking}
\domainid{OA}
\grade{1}
\cluster{Understand and apply properties of operations and the
  relationship between addition and subtraction.}
\standardnumber{3}
\letter{}
\adv{}
\mod{}
\text{Apply properties of operations as strategies to add and
  subtract.\footnote{Students need not use formal terms for these
    properties.}  \clarification{Examples: If $8 + 3 = 11$ is known, then $3 + 8 = 11$ is also
  known. (Commutative property of addition.) To add $2 + 6 + 4$, the
  second two numbers can be added to make a ten, so $2 + 6 + 4 = 2 + 10
  = 12$. (Associative property of addition.)} }

%%% Local Variables: 
%%% mode: latex
%%% TeX-master: "allstandards"
%%% End: 
\standardid{1.OA.4}
\document{Math}
\sectionid{Math K-8}
\anchor{0}
\theme{}
\domain{Operations and Algebraic Thinking}
\domainid{OA}
\grade{1}
\cluster{Understand and apply properties of operations and the
  relationship between addition and subtraction.}
\standardnumber{4}
\letter{}
\adv{}
\mod{}
\text{Understand subtraction as an unknown-addend problem. \clarification{For
  example, subtract $10 - 8$ by finding the number that makes 10 when
  added to 8.}}

%%% Local Variables: 
%%% mode: latex
%%% TeX-master: "allstandards"
%%% End: 
\standardid{1.OA.5}
\document{Math}
\sectionid{Math K-8}
\anchor{0}
\theme{}
\domain{Operations and Algebraic Thinking}
\domainid{OA}
\grade{1}
\cluster{Add and subtract within 20.}
\standardnumber{5}
\letter{}
\adv{}
\mod{}
\text{Relate counting to addition and subtraction (e.g., by counting
  on 2 to add 2).}

%%% Local Variables: 
%%% mode: latex
%%% TeX-master: "allstandards"
%%% End: 
\standardid{1.OA.6}
\document{Math}
\sectionid{Math K-8}
\anchor{0}
\theme{}
\domain{Operations and Algebraic Thinking}
\domainid{OA}
\grade{1}
\cluster{Add and subtract within 20.}
\standardnumber{6}
\letter{}
\adv{}
\mod{}
\text{Add and subtract within 20, demonstrating fluency for addition
  and subtraction within 10. Use strategies such as counting on;
  making ten (e.g., $8 + 6 = 8 + 2 + 4 = 10 + 4 = 14$); decomposing a
  number leading to a ten (e.g., $13 - 4 = 13 - 3 - 1 = 10 - 1 = 9$);
  using the relationship between addition and subtraction (e.g.,
  knowing that $8 + 4 = 12$, one knows $12 - 8 = 4$); and creating
  equivalent but easier or known sums (e.g., adding $6 + 7$ by creating
  the known equivalent $6 + 6 + 1 = 12 + 1 = 13$).}

%%% Local Variables: 
%%% mode: latex
%%% TeX-master: "allstandards"
%%% End: 
\standardid{1.OA.7}
\document{Math}
\sectionid{Math K-8}
\anchor{0}
\theme{}
\domain{Operations and Algebraic Thinking}
\domainid{OA}
\grade{1}
\cluster{Work with addition and subtraction equations.}
\standardnumber{7}
\letter{}
\adv{}
\mod{}
\text{Understand the meaning of the equal sign, and determine if
  equations involving addition and subtraction are true or false. \clarification{For
  example, which of the following equations are true and which are
  false? $6 = 6$, $7 = 8 - 1$, $5 + 2 = 2 + 5$, $4 + 1 = 5 + 2$.}}

%%% Local Variables: 
%%% mode: latex
%%% TeX-master: "allstandards"
%%% End: 
\standardid{1.OA.8}
\document{Math}
\sectionid{Math K-8}
\anchor{0}
\theme{}
\domain{Operations and Algebraic Thinking}
\domainid{OA}
\grade{1}
\cluster{Work with addition and subtraction equations.}
\standardnumber{8}
\letter{}
\adv{}
\mod{}
\text{Determine the unknown whole number in an addition or
  subtraction equation relating three whole numbers. \clarification{For example,
  determine the unknown number that makes the equation true in each of
  the equations $8 + ? = 11$, $5 = \boxvoid - 3$, $6 + 6 = \boxvoid$.}}

%%% Local Variables: 
%%% mode: latex
%%% TeX-master: "allstandards"
%%% End: 
\standardid{1.NBT.1}
\document{Math}
\sectionid{Math K-8}
\anchor{0}
\theme{}
\domain{Number and Operations in Base Ten}
\domainid{NBT}
\grade{1}
\cluster{Extend the counting sequence.}
\standardnumber{1}
\letter{}
\adv{}
\mod{}
\text{Count to 120, starting at any number less than 120. In this
  range, read and write numerals and represent a number of objects
  with a written numeral.}

%%% Local Variables: 
%%% mode: latex
%%% TeX-master: "allstandards"
%%% End: 
\standardid{1.NBT.2}
\document{Math}
\sectionid{Math K-8}
\anchor{0}
\theme{}
\domain{Number and Operations in Base Ten}
\domainid{NBT}
\grade{1}
\cluster{Understand place value.}
\standardnumber{2}
\letter{}
\adv{}
\mod{}
\text{Understand that the two digits of a two-digit number represent
  amounts of tens and ones. Understand the following as special
  cases:}

%%% Local Variables: 
%%% mode: latex
%%% TeX-master: "allstandards"
%%% End: 
\standardid{1.NBT.2a}
\document{Math}
\sectionid{Math K-8}
\anchor{0}
\theme{}
\domain{Number and Operations in Base Ten}
\domainid{NBT}
\grade{1}
\cluster{Understand place value.}
\standardnumber{2}
\letter{a}
\adv{}
\mod{}
\text{10 can be thought of as a bundle of ten ones---called a ``ten.''}

%%% Local Variables: 
%%% mode: latex
%%% TeX-master: "allstandards"
%%% End: 
\standardid{1.NBT.2b}
\document{Math}
\sectionid{Math K-8}
\anchor{0}
\theme{}
\domain{Number and Operations in Base Ten}
\domainid{NBT}
\grade{1}
\cluster{Understand place value.}
\standardnumber{2}
\letter{b}
\adv{}
\mod{}
\text{The numbers from 11 to 19 are composed of a ten and one, two,
  three, four, five, six, seven, eight, or nine ones.}

%%% Local Variables: 
%%% mode: latex
%%% TeX-master: "allstandards"
%%% End: 
\standardid{1.NBT.2c}
\document{Math}
\sectionid{Math K-8}
\anchor{0}
\theme{}
\domain{Number and Operations in Base Ten}
\domainid{NBT}
\grade{1}
\cluster{Understand place value.}
\standardnumber{2}
\letter{c}
\adv{}
\mod{}
\text{The numbers 10, 20, 30, 40, 50, 60, 70, 80, 90 refer to one,
  two, three, four, five, six, seven, eight, or nine tens (and 0
  ones).}

%%% Local Variables: 
%%% mode: latex
%%% TeX-master: "allstandards"
%%% End: 
\standardid{1.NBT.3}
\document{Math}
\sectionid{Math K-8}
\anchor{0}
\theme{}
\domain{Number and Operations in Base Ten}
\domainid{NBT}
\grade{1}
\cluster{Understand place value.}
\standardnumber{3}
\letter{}
\adv{}
\mod{}
\text{Compare two two-digit numbers based on meanings of the tens and
  ones digits, recording the results of comparisons with the symbols
  $>$, $=$, and $<$.}

%%% Local Variables: 
%%% mode: latex
%%% TeX-master: "allstandards"
%%% End: 
\standardid{1.NBT.4}
\document{Math}
\sectionid{Math K-8}
\anchor{0}
\theme{}
\domain{Number and Operations in Base Ten}
\domainid{NBT}
\grade{1}
\cluster{Use place value understanding and properties of operations to
  add and subtract.}
\standardnumber{4}
\letter{}
\adv{}
\mod{}
\text{Add within 100, including adding a two-digit number and a
  one-digit number, and adding a two-digit number and a multiple of
  10, using concrete models or drawings and strategies based on place
  value, properties of operations, and/or the relationship between
  addition and subtraction; relate the strategy to a written method
  and explain the reasoning used. Understand that in adding two-digit
  numbers, one adds tens and tens, ones and ones; and sometimes it is
  necessary to compose a ten.}

%%% Local Variables: 
%%% mode: latex
%%% TeX-master: "allstandards"
%%% End: 
\standardid{1.NBT.5}
\document{Math}
\sectionid{Math K-8}
\anchor{0}
\theme{}
\domain{Number and Operations in Base Ten}
\domainid{NBT}
\grade{1}
\cluster{Use place value understanding and properties of operations to
  add and subtract.}
\standardnumber{5}
\letter{}
\adv{}
\mod{}
\text{Given a two-digit number, mentally find 10 more or 10 less than
  the number, without having to count; explain the reasoning used.}

%%% Local Variables: 
%%% mode: latex
%%% TeX-master: "allstandards"
%%% End: 
\standardid{1.NBT.6}
\document{Math}
\sectionid{Math K-8}
\anchor{0}
\theme{}
\domain{Number and Operations in Base Ten}
\domainid{NBT}
\grade{1}
\cluster{Use place value understanding and properties of operations to
  add and subtract.}
\standardnumber{6}
\letter{}
\adv{}
\mod{}
\text{Subtract multiples of 10 in the range 10--90 from multiples of
  10 in the range 10--90 (positive or zero differences), using concrete
  models or drawings and strategies based on place value, properties
  of operations, and/or the relationship between addition and
  subtraction; relate the strategy to a written method and explain the
  reasoning used.}

%%% Local Variables: 
%%% mode: latex
%%% TeX-master: "allstandards"
%%% End: 
\standardid{2.OA.1}
\document{Math}
\sectionid{Math K-8}
\anchor{0}
\theme{}
\domain{Operations and Algebraic Thinking}
\domainid{OA}
\grade{2}
\cluster{Represent and solve problems involving addition and
  subtraction.}
\standardnumber{1}
\letter{}
\adv{}
\mod{}
\text{Use addition and subtraction within 100 to solve one- and
  two-step word problems involving situations of adding to, taking
  from, putting together, taking apart, and comparing, with unknowns
  in all positions, e.g., by using drawings and equations with a
  symbol for the unknown number to represent the problem.\footnote{See
  Glossary, Table 1.}}

%%% Local Variables: 
%%% mode: latex
%%% TeX-master: "allstandards"
%%% End: 
\standardid{2.OA.2}
\document{Math}
\sectionid{Math K-8}
\anchor{0}
\theme{}
\domain{Operations and Algebraic Thinking}
\domainid{OA}
\grade{2}
\cluster{Add and subtract within 20.}
\standardnumber{2}
\letter{}
\adv{}
\mod{}
\text{Fluently add and subtract within 20 using mental
  strategies.\footnote{See standard 1.OA.6 for a list of mental
    strategies.} 
  By end of Grade 2, know from memory all sums of two one-digit
  numbers.}

%%% Local Variables: 
%%% mode: latex
%%% TeX-master: "allstandards"
%%% End: 
\standardid{2.OA.3}
\document{Math}
\sectionid{Math K-8}
\anchor{0}
\theme{}
\domain{Operations and Algebraic Thinking}
\domainid{OA}
\grade{2}
\cluster{Work with equal groups of objects to gain foundations for
  multiplication.}
\standardnumber{3}
\letter{}
\adv{}
\mod{}
\text{Determine whether a group of objects (up to 20) has an odd or
  even number of members, e.g., by pairing objects or counting them by
  2s; write an equation to express an even number as a sum of two
  equal addends.}

%%% Local Variables: 
%%% mode: latex
%%% TeX-master: "allstandards"
%%% End: 
\standardid{2.OA.4}
\document{Math}
\sectionid{Math K-8}
\anchor{0}
\theme{}
\domain{Operations and Algebraic Thinking}
\domainid{OA}
\grade{2}
\cluster{Work with equal groups of objects to gain foundations for
  multiplication.}
\standardnumber{4}
\letter{}
\adv{}
\mod{}
\text{Use addition to find the total number of objects arranged in
  rectangular arrays with up to 5 rows and up to 5 columns; write an
  equation to express the total as a sum of equal addends.}

%%% Local Variables: 
%%% mode: latex
%%% TeX-master: "allstandards"
%%% End: 
\standardid{2.NBT.1}
\document{Math}
\sectionid{Math K-8}
\anchor{0}
\theme{}
\domain{Number and Operations in Base Ten}
\domainid{NBT}
\grade{2}
\cluster{Understand place value.}
\standardnumber{1}
\letter{}
\adv{}
\mod{}
\text{Understand that the three digits of a three-digit number
  represent amounts of hundreds, tens, and ones; e.g., 706 equals 7
  hundreds, 0 tens, and 6 ones. Understand the following as special
  cases:}

%%% Local Variables: 
%%% mode: latex
%%% TeX-master: "allstandards"
%%% End: 
\standardid{2.NBT.1a}
\document{Math}
\sectionid{Math K-8}
\anchor{0}
\theme{}
\domain{Number and Operations in Base Ten}
\domainid{NBT}
\grade{2}
\cluster{Understand place value.}
\standardnumber{1}
\letter{a}
\adv{}
\mod{}
\text{100 can be thought of as a bundle of ten tens---called a
  ``hundred.''}

%%% Local Variables: 
%%% mode: latex
%%% TeX-master: "allstandards"
%%% End: 
\standardid{2.NBT.1b}
\document{Math}
\sectionid{Math K-8}
\anchor{0}
\theme{}
\domain{Number and Operations in Base Ten}
\domainid{NBT}
\grade{2}
\cluster{Understand place value.}
\standardnumber{1}
\letter{b}
\adv{}
\mod{}
\text{The numbers 100, 200, 300, 400, 500, 600, 700, 800, 900 refer
  to one, two, three, four, five, six, seven, eight, or nine hundreds
  (and 0 tens and 0 ones).}

%%% Local Variables: 
%%% mode: latex
%%% TeX-master: "allstandards"
%%% End: 
\standardid{2.NBT.2}
\document{Math}
\sectionid{Math K-8}
\anchor{0}
\theme{}
\domain{Number and Operations in Base Ten}
\domainid{NBT}
\grade{2}
\cluster{Understand place value.}
\standardnumber{2}
\letter{}
\adv{}
\mod{}
\text{Count within 1000; skip-count by 5s, 10s, and 100s.}

%%% Local Variables: 
%%% mode: latex
%%% TeX-master: "allstandards"
%%% End: 
\standardid{2.NBT.3}
\document{Math}
\sectionid{Math K-8}
\anchor{0}
\theme{}
\domain{Number and Operations in Base Ten}
\domainid{NBT}
\grade{2}
\cluster{Understand place value.}
\standardnumber{3}
\letter{}
\adv{}
\mod{}
\text{Read and write numbers to 1000 using base-ten numerals, number
  names, and expanded form.}

%%% Local Variables: 
%%% mode: latex
%%% TeX-master: "allstandards"
%%% End: 
\standardid{2.NBT.4}
\document{Math}
\sectionid{Math K-8}
\anchor{0}
\theme{}
\domain{Number and Operations in Base Ten}
\domainid{NBT}
\grade{2}
\cluster{Understand place value.}
\standardnumber{4}
\letter{}
\adv{}
\mod{}
\text{Compare two three-digit numbers based on meanings of the
  hundreds, tens, and ones digits, using $>$, =, and $<$ symbols to record
  the results of comparisons.}

%%% Local Variables: 
%%% mode: latex
%%% TeX-master: "allstandards"
%%% End: 
\standardid{2.NBT.5}
\document{Math}
\sectionid{Math K-8}
\anchor{0}
\theme{}
\domain{Number and Operations in Base Ten}
\domainid{NBT}
\grade{2}
\cluster{Use place value understanding and properties of operations to
  add and subtract.}
\standardnumber{5}
\letter{}
\adv{}
\mod{}
\text{Fluently add and subtract within 100 using strategies based on
  place value, properties of operations, and/or the relationship
  between addition and subtraction.}

%%% Local Variables: 
%%% mode: latex
%%% TeX-master: "allstandards"
%%% End: 
\standardid{2.NBT.6}
\document{Math}
\sectionid{Math K-8}
\anchor{0}
\theme{}
\domain{Number and Operations in Base Ten}
\domainid{NBT}
\grade{2}
\cluster{Use place value understanding and properties of operations to
  add and subtract.}
\standardnumber{6}
\letter{}
\adv{}
\mod{}
\text{Add up to four two-digit numbers using strategies based on place
  value and properties of operations.}

%%% Local Variables: 
%%% mode: latex
%%% TeX-master: "allstandards"
%%% End: 
\standardid{2.NBT.7}
\document{Math}
\sectionid{Math K-8}
\anchor{0}
\theme{}
\domain{Number and Operations in Base Ten}
\domainid{NBT}
\grade{2}
\cluster{Use place value understanding and properties of operations to
  add and subtract.}
\standardnumber{7}
\letter{}
\adv{}
\mod{}
\text{Add and subtract within 1000, using concrete models or drawings
  and strategies based on place value, properties of operations,
  and/or the relationship between addition and subtraction; relate the
  strategy to a written method. Understand that in adding or
  subtracting three-digit numbers, one adds or subtracts hundreds and
  hundreds, tens and tens, ones and ones; and sometimes it is
  necessary to compose or decompose tens or hundreds.}

%%% Local Variables: 
%%% mode: latex
%%% TeX-master: "allstandards"
%%% End: 
\standardid{2.NBT.8}
\document{Math}
\sectionid{Math K-8}
\anchor{0}
\theme{}
\domain{Number and Operations in Base Ten}
\domainid{NBT}
\grade{2}
\cluster{Use place value understanding and properties of operations to
  add and subtract.}
\standardnumber{8}
\letter{}
\adv{}
\mod{}
\text{Mentally add 10 or 100 to a given number 100--900, and mentally
  subtract 10 or 100 from a given number 100--900.}

%%% Local Variables: 
%%% mode: latex
%%% TeX-master: "allstandards"
%%% End: 
\standardid{2.NBT.9}
\document{Math}
\sectionid{Math K-8}
\anchor{0}
\theme{}
\domain{Number and Operations in Base Ten}
\domainid{NBT}
\grade{2}
\cluster{Use place value understanding and properties of operations to
  add and subtract.}
\standardnumber{9}
\letter{}
\adv{}
\mod{}
\text{Explain why addition and subtraction strategies work, using
  place value and the properties of operations.\footnote{Explanations
    may be supported by drawings or objects.}}

%%% Local Variables: 
%%% mode: latex
%%% TeX-master: "allstandards"
%%% End: 
